\documentclass[../main.tex]{subfiles}
\graphicspath{{\subfix{../img/}}}
\begin{document}
	\subsection{Ricapitolando}
	Neila scorsa sezione abbiamo introdotto i semiconduttori e ne abbiamo evidenziato le principali differenze rispetto ai conduttori.\\\\
	Abbiamo inoltre modellizzato i semiconduttori come:
	$$\overline{J} = \overline{J}_n + \overline{J}_p$$
	dove $\overline{J}_n$ è la densità di corrente dovuta agli elettroni e $\overline{J}_p$ è la densità di corrente dovuta alle lacune. \\
	Nel casobidimensionale (che è il caso di nostro interesse), tali componenti si possono esprimere come:
	$$J_n = -q\mu_n n E + q D_n \frac{dn}{dx}$$
	$$J_p = q\mu_p p E - q D_p \frac{dp}{dx}$$
	Infine abbiamo sottolineato come la mobilità elettronica e la mobilità delle lacune possano essere scritte come:
	$$D_n = \frac{KT}{q} \mu_n$$
	$$D_p = \frac{KT}{q} \mu_p$$
	con $D_n$ tipicamente più grande di due o tre volte rispetto a $D_p$.
	\subsection{Applicazione del modello}
	\label{sub:app-mod}
	\begin{center}
		\begin{figure}[h]
			\centering
			\includegraphics[scale=0.55]{7.png}
			\caption{Semiconduttore con drogaggi differenti all’equilibrio ($T = 0^\degree K$)}
			\label{fig:im-7}
		\end{figure}
	\end{center}
	In un materiale di tipo $n$ dorgato non uniformemente abbiamo:
	$$J_n = 0 \implies q\mu_n n E + q D_n \frac{dn}{dx} = 0$$
	ma sappiamo che è:
	$$E = -\frac{d\varphi}{dx}$$
	$$D_n = \frac{KT}{q}\mu_n$$
	e quindi, sostituendo, possiamo scrivere:
	$$0=-\cancel{q}\cancel{\mu_n} n \frac{d\varphi}{dx} + \cancel{q} \frac{KT}{q}\cancel{\mu_n} \frac{dn}{dx} \implies \frac{d\varphi}{dx} = \frac{KT}{q} \cdot \frac{1}{n} \cdot \frac{dn}{dx}$$
	$$\rightarrow \int_{\varphi(n_1)=\varphi}^{\varphi(n_1)=\varphi_2} \frac{d \varphi}{\cancel{dx}} = \int_{n(x_1)=n_1}^{n(x_2)=n_1} \frac{KT}{q} \frac{1}{n} \frac{dn}{\cancel{dx}} = \frac{KT}{q} \int_{n_1}^{n_2} \frac{1}{n} dn$$
	che, integrando, porta a:
	$$\varphi_2 - \varphi = \frac{KT}{q} \cdot \ln \frac{n_2}{n_1}$$
	Allora:
	\begin{itemize}
		\item Conoscendo la differenza di concentrazione in due punti, possiamo calcolare il potenziale
		\item Conoscendo la differenza di potenziale in due punti, possiamo calcolare la concentrazione
	\end{itemize}
	Da questa considerazione possiamo decidere arbitrariamente che esiste un punto $x_0$ in cui, all’equilibrio, sia:
	$$\begin{cases} \varphi = 0 \\ n = 0 \\ p = 0 \end{cases}$$
	Allora abbiamo:
	$$\varphi(x) = \frac{KT}{q} \cdot \ln \frac{n(x)}{n_0} \iff n(x) = n_0 \cdot e^{\frac{q\varphi(x)}{KT}}$$
	Lo stesso ragionamento si può applicare imponendo $J_p = 0$. \\
	In questo caso sarebbe:
	$$\varphi(x) = -\frac{KT}{q} \cdot \ln \frac{p(x)}{p_0} \iff p(x) = p_0 \cdot e^{-\frac{q\varphi(x)}{KT}}$$
	\subsection{La Giunzone \textit{pn}}
	\label{sub:pn}
	Premessa: \textbf{La giunzione $pn$ è il primo dispositivo elettronico che verrà introdotto in questo corso, la sua struttura sarà alla base di gran parte dei dispositivi elettronici più complessi}.\\\\
	Mantenendoci nelle condizioni di equilibrio e di bidimensionalità, vogliamo drograre un blocco di semiconduttore con concentrazioni uniformi di atomi accettori e donatori.\\\\
	Il punto in cui si passa dalla regione drogata di tipo $p$ (cioè con atomi accettori) e la regione di tipo $n$ (con atomi donatori) è detto giunzione metallurgica o più semplicemente giunzione.\\\\
	In condizione di equilibrio sappiamo che è:
	$$J_n = J_p = 0$$
	Possiamo aspettarci che la maggior parte dei fenomeni di diffusione si concentri nei pressi del punto di giunzione.\\
	Conseguentemente possiamo anche ipotizzare che esista una coordinata $w_n$ ed una sua duale $-w_p$ oltre le quali non si notano gli effetti perturbativi dovuti alla giunzione. \\
	Allora il dispositivo è diviso in quattro regioni:
	\begin{enumerate}
		\item una regione drogata $p$ che non risente della giunzione
		\item una regione drogata ( $p$ ) che risente della giunzione
		\item una regione drogata ($n$) che risente della giunzione
		\item una regione drogata ($n$) che non risente della giunzione
	\end{enumerate}
	I casi 1 e 4 sono già stati studiati in precedenza.
	\begin{center}
		\begin{figure}[h]
			\centering
			\includegraphics[scale=0.55]{8.png}
			\caption{Schema di una giunzione $pn$ e delle quattro zone in cui si può ipotizzare di suddividerla}
			\label{fig:im-8}
		\end{figure}
	\end{center}
	\subsubsection{Intervallo $(-\infty; -w_p)$}
	\label{subsub:int-pn1}
	Questa regione è detta zona neutra di tipo $p$.\\\\
	In questo intervallo abbiamo:
	$$p \approx \underbrace{N_A}_{\text{N di accettatori}}$$
	$$n \approx \frac{\overbrace{n_i^2}^{\text{Ricordando }p\cdot n = n_i^2}}{N_A}$$
	$$\rho = q(N_D - N_A + p - n) = 0$$
	Inoltre deve essere valida l’uguaglianza:
	$$J_p = q\mu_p p E - qD_p \frac{dp}{dx} = 0$$
	ma siccome il secondo membro è nullo (la derivata di una costante è sempre nulla) deve essere nullo anche il primo membro e dunque sarà: $E = 0$ $(\underbrace{E}_{=0} = \underbrace{- \frac{d \varphi}{dx}}_{=0})$.\\\\
	Il potenziale dovrà essere costante (se il campo elettrico, che è la derivata del potenziale, è nullo il potenziale dovrà essere necessariamente costante). \\
	Allora possiamo decidere che il potenziale $\varphi$ sia nullo.
	\subsubsection{Intervallo $(-w_p; w_n)$}
	\label{subsub:int-pn2}
	Sappiamo con ragionevolezza che nell’intervallo $(-w_p; w_n)$ il potenziale sarà:
	$$0 < \varphi < \psi_b$$
	Tuttavia sappiamo anche che, per $\varphi$ anche di poco maggiore di zero (risultato necessariamente verificato per $x > -w_p$), sarà:
	$$\varphi > 0$$
	$$p(x) = p_0 \cdot e^{-\frac{q\varphi(x)}{KT}}$$
	$$p(x) \gg p_0 \approx N_A$$
	Dalla parte opposta, avremo invece, necessariamente:
	$$\varphi < \psi$$
	$$n(x) = n_0 \cdot e^{-\frac{q\varphi(x)}{kT}}$$
	$$n(x) \ll n_0 \approx N_D$$
	Dunque in tutto l’intervallo $(-w_p; w_n)$ \textbf{la concentrazione di elettroni e lacune è trascurabile}. \\
	Questa regione è detta \textbf{regione svuotata}. \\
	Il termine “svuotata” non indica che non vi siano portatori di carica ma solo che essi sono in una concentrazione trascurabile rispetto alle concentrazioni delle regioni neutre.
	\subsubsection{Intervallo $(-w_p; 0)$}
	\label{subsub:int-pn3}
	In questo intervallo non abbiamo portatori mobili. ({\color{red}* DA RIVEDERE})\\\\
	Allora in questa regione la densità di carica è:
	$$\rho = q(-N_A + p - n) = -qN_A \qquad \{ p,n \to 0$$
	poiché $p$ ed $n$ sono trascurabili per piccole variazioni della concentrazione.\\\\
	Inoltre:
	$$\frac{dE}{dx} = \frac{\rho}{\varepsilon}$$
	Conseguentemente:
	$$\frac{dE}{dx} = \frac{-qN_A}{\varepsilon}$$
	da cui:
	$$dE = -\frac{qN_A}{\varepsilon}dx \overbrace{\implies}^{\int_{E(-w_p)}^{E(n)} dE = \int_{-w_p}^{n} - \frac{qE}{\varepsilon} dn} E(x) - \underbrace{E(-w_p)}_{=0} = -\frac{qN_A}{\varepsilon}(x + w_p)$$
	Per continuità con l’intervallo $(-\infty; -w_p)$, avremo $E(-w_p) = 0$.
	$$E(0) = -\frac{qN_A}{\varepsilon}(w_p)$$
	Allora nell’intervallo in questione il campo elettrico è una retta con coefficiente angolare negativo.\\\\
	Infine il potenziale è:
	$$E = - \frac{d \varphi}{dx} \rightarrow \frac{d\varphi}{dx} = \frac{qN_A}{\varepsilon}(x + w_p) \implies d\varphi = \frac{qN_A}{\varepsilon}(x + w_p)dx$$
	che, integrando da $-w_p$ ad $x$, porta a:
	$$\varphi(x) - \varphi(w_p) = \frac{qN_A}{\varepsilon} \cdot \frac{(x + w_p)^2}{2} = \frac{qN_A}{2\varepsilon}(x - w_p)^2$$
	Anche in questo caso, per continuità con l’intervallo $(-\infty; -w_p)$, avremo $\varphi(w_p) = 0$.
	$$\varphi(0) = \frac{qN_A}{2\varepsilon}(w_p)^2$$
	\subsubsection{Intervallo $(0; w_n)$}
	\label{subsub:int-pn4}
	Applicando ragionamenti similari a quelli visti nel,'intervallo \ref{subsub:int-pn3} otteniamo una densità di carica:
	$$\rho = qN_D$$
	La positività della costante è dovuta al fatto che, in questa regione, gli ioni fissi portano carica positiva.\\\\
	Per il campo elettrico avremo invece:
	$$E\left(x\right)\overbrace{=}^{\frac{dE}{dx}=\frac{qN_{0}}{E}\rightarrow\int_{E(w_n)}^{E(n)} dE=\int_{w_n}^{n}dx}\frac{qN_D}{\varepsilon}\left(x-w_n\right)$$
	che è una retta con coefficiente angolare positivo ed ordinata di origine negativa.
	$$E\left(0\right)=\frac{qN_D}{\varepsilon}\left(-w_n\right)$$
	Infine il potenziale si calcola in modo simile a quanto visto per l’intervallo $(-w_p; 0)$.
	$$\begin{cases} E=\frac{qN_d}{\varepsilon}\left(x-w_n\right) \\ E=-\frac{d\varphi}{dx}\end{cases}
	\;
	\Rightarrow-d\varphi=\frac{qN_d}{\varepsilon}\left(x-w_n\right)dx
	\rightarrow \int_{\varphi(w_{n})}^{\varphi(0)} -d\varphi \int_{w_{n}}^{0}\frac{q N_D}{\varepsilon}\left(x-w_{n}\right)dx$$
	Integrando da $\varphi\left(0\right)$ a $\varphi\left(w_n\right)=\psi_b$ otteniamo:
	$$\varphi\left(x\right)=\psi_b-\frac{qN_D}{2\varepsilon}\left(x-w_p\right)^2$$
	\subsubsection{Intervallo $(w_n; +\infty)$}
	\label{subsub:int-pn5}
	Questa regione è detta \textbf{zona neutra di tipo $n$}.\\\\
	In questa regione valgono tutti i ragionamenti visti nell'intervallo \ref{subsub:int-pn1} con l’unica differenza che non possiamo imporre arbitrariamente il valore di $\varphi$.\\\\
	In questo intervallo abbiamo:
	$$n\approx N_D$$
	$$p\approx\frac{n_i^2}{N_D}$$
	$$\rho=q\left(N_D-N_A+p-n\right)=0$$
	Inoltre deve essere valida l’uguaglianza:
	$$J_n=q\mu_{np}pE-qD_n\frac{dn}{dx}=0\implies E=0$$
	In questo caso non possiamo porre nuovamente il potenziale a zero poiché esso deve essere misurato in relazione al valore $\varphi=0$ imposto in precedenza. \\
	Dalla relazione ricavata nella sezione \ref{sub:app-mod} abbiamo:
	$$\varphi\left(x>w_n\right)=\frac{KT}{q}\cdot\ln\frac{N_D N_A}{n_i^2}$$
	e, siccome $\frac{KT}{q}$ \textbf{deve essere un potenziale} solitamente si indica con $V_T$ ed è detta \textbf{tensione termica}.\\
	A temperatura ambiente abbiamo:
	$$V_T=25 \ \mathrm{mV}$$
	Il valore del potenziale nella zona neutra $n$ è indicato con $\psi_b$ ed è detto \textbf{potenziale barriera} poiché deve fungere da barriera contro la diffusione dei portatori di carica.
	\subsubsection{Grafici}
	I grafici seguenti (grafici \ref{fig:image3} e grafici \ref{fig:image4}) mostrano le principali caratteristiche di una giunzione $p_n$ nelle quattro regioni.
	\begin{figure}[h]
		
		\begin{subfigure}{0.5\textwidth}
			\includegraphics[width=0.9\linewidth]{9.png} 
			\caption{Concentrazione di atomi accettori e donatori in un diodo a giunzione $pn$ a temperatura ambiente}
			\label{fig:im-9}
		\end{subfigure}
		\begin{subfigure}{0.5\textwidth}
			\includegraphics[width=0.9\linewidth]{10.png}
			\caption{Grafico della densità di carica in un diodo a giunzione $pn$}
			\label{fig:im-10}
		\end{subfigure}
		
		\caption{}
		\label{fig:image3}
	\end{figure}
	\begin{figure}[h]
		
		\begin{subfigure}{0.5\textwidth}
			\includegraphics[width=0.9\linewidth]{11.png} 
			\caption{Campo elettrico in un diodo a giunzione $pn$}
			\label{fig:im-11}
		\end{subfigure}
		\begin{subfigure}{0.5\textwidth}
			\includegraphics[width=0.9\linewidth]{12.png}
			\caption{Potenziale in un diodo a giunzione $pn$}
			\label{fig:im-12}
		\end{subfigure}
		
		\caption{}
		\label{fig:image4}
	\end{figure}
	Si noti che è:
	$$q N _ { A } w _ { p } = q N _ { D } w _ { n }$$
	Da questa relazione notiamo che la regione si svuota in maniera inversamente proporzionale alla concentrazione poiché le aree azzurre in figura \ref{fig:im-10} devono essere uguali.\\\\
	Oltre ad indicare che la giunzione si comporta quasi come un condensatore, questa relazione ci consente di ridurre di una ingognita le nostre equazioni.
	\subsection{Validazione}
	I risultati ed i grafici ottenuti nella sezione \ref{sub:pn}, si basano su una ipotesi: che $-w_p$ e $w_n$ siano entrambi compatibili con le dimensioni fisiche del dispositivo. \\
	Dobbiamo ora valutare se l’ipotesi è corretta o meno.\\\\
	Sappiamo che il campo elettrico si può calcolare in due modi differenti a seconda che consideriamo la zona svuotata $n$ o nella zona svuotata $p$. \\
	Siccome le due zone si congiungono in $x = 0$, in quel punto il campo dovrà, per continuità, essere lo stesso sia che venga calcolato da destra sia che venga calcolato da destra. \\
	Abbiamo ricavato in precedenza che deve essere:
	$$qN_A w_p = qN_D w_n$$
	Ragionando in modo analogo per il potenziale otteniamo che deve risultare:
	$$\varphi(x)=\psi_{b}-\frac{qN_{A}}{2\varepsilon}\left(x-w_{P}\right)^{2} \overbrace{\rightarrow}^{\text{Le zone si congiungono in }x=0} \frac{qN_A}{2\epsilon} w_p^2 = \psi_b - \frac{qN_D}{2\epsilon} w_n^2$$
	che è un’equazione in due incognite.\\\\
	Allora possiamo mettere le due equazioni a sistema e risolvere per sostituzione.
	$$\begin{cases} 
			qN_A w_p = qN_D w_n \\ 
			\frac{qN_A}{2\varepsilon} w_p^2 = \psi_b - \frac{qN_D}{2\varepsilon} w_n^2 
		\end{cases} 
		\implies 
		\begin{cases} 
			w_p = \frac{N_D}{N_A} w_n \\ 
			\psi_b = \frac{q}{2\varepsilon} \left( N_A w_p^2 + N_A w_n^2 \right) 
		\end{cases}$$
	$$\psi_b = \frac{q}{2\varepsilon} \left( N_A \left( \frac{N_D}{N_A} w_n \right)^2 + N_A w_n^2 \right) = \frac{qN_D^2 w_n^2}{2\varepsilon} \left( \frac{1}{N_A} + \frac{1}{N_D} \right)$$
	Sfruttando ora il fatto che
	$$\psi_b = \frac{KT}{q} \ln \frac{N_D N_A}{n_i^2}$$
	e sciogliendo il quadrato facendo la radice di ambo i membri, ricaviamo un valore numerico per $w_n$ e $w_p$:
	$$w_n = \frac{1}{N_D} \sqrt{\frac{2\varepsilon \psi_b}{q \left( \frac{1}{N_A} + \frac{1}{N_D} \right)}}$$
	$$w_p = \frac{1}{N_A} \sqrt{\frac{2\varepsilon \psi_b}{q \left( \frac{1}{N_A} + \frac{1}{N_D} \right)}}$$
	Tali valori, con parametri tipici, si attestano sulla dimensione di qualche micron. \\
	Allora \textbf{è sufficiente che il blocchetto di semiconduttore abbia dimensioni superiori a qualche micron} perché l’ipotesi sia soddisfatta ed il ragionamento sia valido.
	\subsection{Considerazioni}
	Possiamo inoltre calcolare l’ampiezza totale della regione svuotata.
	$$w \doteq w_p + w_n = \left( \frac{1}{N_A} + \frac{1}{N_D} \right) \sqrt{\frac{2\varepsilon\psi_b}{q \left( \frac{1}{N_A} + \frac{1}{N_D} \right)}} = \sqrt{\frac{2\varepsilon\psi_b}{q} \left( \frac{1}{N_A} + \frac{1}{N_D} \right)}$$
	Da questo risultato notiamo che la regione svuotata è tanto più piccola quanto più è drogato il materiale. \\
	Conseguentemente, se una parte del materiale è più drogata dell’altra, avremo una regione svuotata asimmettrica più piccola nella zona maggiormente drogata.\\\\
	Inoltre, dai grafici, notiamo che il campo elettrico ha verso opposto rispetto all’asse $x$. \\
	Dunque il campo ha un effetto di trascinamento delle lacune verso la zona $p$ e degli elettroni verso la zona $n$ che contrasta l’effetto diffusivo dovuto alla differenza di concentrazione. \\
	Abbiamo dunque una differenza di potenziale tra le due zone.\\\\
	Questa differenza di potenziale non produce tuttavia energia all’esterno (si avrebbe una violazione del principio di conservazione dell’energia): attaccando una lampadina ai capi delle regioni drogate questa non si accende. \\
	Questo risultato si spiega considerando che la giunzione tra il metallo ed il semiconduttore ha effetti simili alla giunzione $pn$ e deve essere nuovamente studiata in dettaglio. \\
	Anche fra la zona $p$ ed il metallo dunque si sviluppa una differenza di potenziale (detta \textbf{differenza di potenziale di contatto} ed indicata da $\psi_{cp}$ o $\psi_{cn}$) (Dove $C$ indica la zona conduttiva). \\
	Allora dovrà essere:
	$$\underbrace{\psi_b}_{\underset{\text{di Barriera}}{\text{Potenziale}}} + \underbrace{\psi_{cp}}_{\underset{\text{Contatto Zona $p$ e $c$}}{\text{Potenziale di}}} + \underbrace{\psi_{cn}}_{\underset{\text{Contatto Zona $c$ e $n$}}{\text{Potenziale di}}} = 0$$
	La differenza di potenziale di contatto è dipendente solamente dalla configurazione fisica del sistema e non dalla corrente che lo attraversa. \\
	Un contatto di questo tipi è detto \textbf{contatto ohmico}.
	\begin{center}
		\begin{figure}[h]
			\centering
			\includegraphics[scale=0.55]{13.png}
			\caption{Diodo a giunzione $pn$ collegato ad un filo di rame}
			\label{fig:im-13}
		\end{figure}
	\end{center}
	\subsection{Funzionamento}
	\label{sub:funz-pn}
	Fino ad ora abbiamo studiato il comportamento del nostro oggetto in condizione di equilibrio. \\
	Vogliamo adesso studiarne il funzionamento fuori equilibrio: agganciamo un generatore di tesione $V$ al sistema (figura \ref{fig:im-14}).
	\begin{center}
		\begin{figure}[h]
			\centering
			\includegraphics[scale=0.55]{14.png}
			\caption{Applicazione di un generatore di tensione al diodo}
			\label{fig:im-14}
		\end{figure}
	\end{center}
	All’equilibrio, cioè con $V = 0$, abbiamo, per la legge di Kirchoff:
	$$\psi_{cp} + \psi_{b_0} + \psi_{cn} = 0$$
	dove $\psi_{b_0}$ è il potenziale di barriera all’equilibrio.\\\\
	Usciamo dalla condizione di equilibrio: al variare di $V$ dovrà variare anche $\psi_b$. \\
	Tuttavia, lo abbiamo studiato in precedenza, le zone neutre si comportano come resistenze. \\
	Le cadute di potenziale sulle zone neutre sono indicate con $\psi_{znp}$ e $\psi_{znn}$. \\
	Allora, nuovamente per la legge di Kirchoff, abbiamo:
	$$V + \psi_{cp} - \psi_{znp} + \psi_b - \psi_{znn} - \psi_{cn} = 0$$
	Ipotizziamo che i contatti siano ohmici e supponiamo che i ragionamenti qualitativi fatti nella condizione di equilibrio siano ancora validi\footnote{Supponiamo cioè che di ragionare in condizione di piccolo scostamento dalla condizione di equilibrio.}. \\
	Sotto queste ipotesi, la corrente sarà qualitativamente piccola e, conseguentemente, le cadute di potenziale sulle zone neutre siano qualitativamente piccole. \\
	Quantifichiamo il termine \textbf{piccolo} come:
	$$\psi_{znp}, \psi_{znn} \ll \psi_b$$
	Allora possiamo semplificare l’equazione di maglia\footnote{Questa semplificazione non sarà più accettabile quando l’ipotesi di piccolo scostamento cadrà.} in:
	$$0 + \psi_{cp} - 0 + \psi_{b_0} - 0 - \psi_{cn} = 0$$
	Sottraendo membro a membro le due equazioni ottenute per la condizione fuori equilibrio abbiamo:
	$$V + \psi_b - \psi_{b_0} = 0$$
	da cui ricaviamo che il potenziale di barriera, in condizione di piccolo scostamento dall’equilibrio, è pari al potenziale di barriera in condizione di equilibrio diminuito della tensione applicata:
	$$\psi_b = \psi_{b_0} - V$$
	Dunque applicando una tensione positiva abbassiamo la barriera di potenziale mentre applicando una tensione negativa alziamo la barriera di potenziale.
	$$V > 0 \Longrightarrow \psi_b > \psi_{b_0}$$
	$$V < 0 \Longrightarrow \psi_b < \psi_{b_0}$$
	Sfruttando il grafico di figura \ref{fig:im-15} notiamo due comportamenti differenti in base alla tensione applicata.
	\begin{center}
		\begin{figure}[h]
			\centering
			\includegraphics[scale=0.55]{15.png}
			\caption{Una differenza di potenziale modifica l’influenza del campo elettrico}
			\label{fig:im-15}
		\end{figure}
	\end{center}
	\quad\\
	Ricordiamo che $E= -grad(V)$.\\\\
	Se $V$ è positivo, abbiamo un abbassamento del campo ed il termine diffusivo è più importante del termine di trascinamento. \\
	Di conseguenza la corrente sarà positiva.\\\\
	Se invece $V$ è negativo abbiamo un innalzamento del campo ed il termine di trascinamento è più forte rispetto al termine di diffusione. \\
	Di conseguenza la corrente sarà negativa.
	\begin{table}[h]
		\centering
		\begin{tabular}{|c|c|c|}
			\hline
			\textbf{Tensione} & \textbf{Effetti} & \textbf{Corrente} \\ \hline
			$V > 0$ & Trascinamento $<$ Diffusione & $I > 0$ \\ \hline
			$V = 0$ & Trascinamento = Diffusione & $I = 0$ \\ \hline
			$V < 0$ & Trascinamento $>$ Diffusione & $I < 0$ \\ \hline
		\end{tabular}
		\caption{Effetto vincente rispetto a differenti valori di tensione}
		\label{tab:table-2}
	\end{table}
	La differenza fondamentale è dovuta alla differenza della popolazione su cui agisce l’effetto:
	\begin{itemize}
		\item La corrente di diffusione agisce sui portatori maggioritari e quindi basta una ridotta variazione di tensione per avere una grande variazione di corrente
		\item La corrente di trascinamento agisce sui portatori minoritari e quindi una ridotta variazione di tensione porta ad una variazione di corrente molto ridotta.
	\end{itemize}
	Il comportamento è graficato in figura \ref{fig:im-16}.
	\begin{center}
		\begin{figure}[h]
			\centering
			\includegraphics[scale=0.55]{16.png}
			\caption{Caratteristica e simbolo grafico di un diodo a giunzione $pn$}
			\label{fig:im-16}
		\end{figure}
	\end{center}
	Il dispositivo appena studiato è detto \textbf{diodo a giunzione $pn$} e viene rappresentato simbolicamente come una freccia che punta nella direzione in cui la corrente procede con facilità ed una barriera per il verso opposto.\\\\
	Il diodo a giunzione $pn$ ha due regioni di funzionamento: una detta \textbf{diretta} (ON) (quando $V$ ed $I$ sono positive) ed una detta \textbf{inversa} (OFF) (con $V$ ed $I$ negative).\\\\
	La caratteristica del diodo a giunzione $pn$ è modellabile come:
	$$I = I_S \left( e^{\frac{qV}{kT}} - 1 \right)$$
	In questo modello:
	\begin{itemize}
		\item Il valore ($-1$) è necessario per far passare dall’origine la curva esponenziale che altrimenti passerebbe nel punto ($I = 1$) per una tensione nulla
		\item La costante ($I_{S}$) (la corrente inversa di saturazione) è il valore di corrente cui tende asintoticamente il valore della corrente quando si applicano valori di corrente molto negativi
	\end{itemize}
	Sperimentalmente si nota che il modello così ricavato funziona bene solo quando tutte le ipotesi iniziali sono rispettate. \\
	Tuttavia se la tensione eccede un certo valore negativo, la corrente ricomincia a crescere a grande velocità risultando fortemente negativa.\\\\
	La tensione per cui avviene questo fenomeno è detta tensione di scarica o \textbf{tensione di break-down} ed è dovuta al fatto che, quando la tensione negativa eccede un certo valore, anche il campo elettrico inizia a generare coppie elettrone lacuna.\\\\
	L’ipotesi caduta è dunque la trascurabilità dell’azione del campo rispetto all’azione diffusiva.\\\\
	Aumentando troppo la tensione positiva e tracciando la curva su scala logaritmica, otteniamo il comportamento graficato in figura \ref{fig:im-17}.
	\begin{center}
		\begin{figure}[h]
			\centering
			\includegraphics[scale=0.55]{17.png}
			\caption{Corrente in funzione della tensione in scala logaritmica}
			\label{fig:im-17}
		\end{figure}
	\end{center}
	L’ipotesi caduta in questo caso è la trascurabilità delle cadute sulle zone neutre poiché non è più vero che sia:
	$$\psi_{znp}, \psi_{znn} \ll \psi_b$$
	\begin{center}
		\begin{figure}[h]
			\centering
			\includegraphics[scale=0.55]{18.png}
			\caption{Esempi di diodi in commercio}
			\label{fig:im-18}
		\end{figure}
	\end{center}
	\subsection{Utilità del Diodo}
	Premessa: nella scorsa sezione: \ref{sub:funz-pn} abbiamo introdotto il diodo a giunzione $pn$. \\
	Abbiamo inoltre ricavato un modello matematico per il diodo nell’ipotesi di piccolo scostamento dall’equilibrio.
	\begin{gather*}
		I = I_S \underbrace{\left( e^{\frac{V}{V_T}} - 1 \right)}_{\text{Risulta essere un volume adimensionale}} \\
		\downarrow \\
		{\color{red}+1}+\frac{I_D}{\color{teal}I_S} = {\color{teal}I_S} - \left( e^{\frac{V_D}{V_T}} {\color{red}-1} \right) \longrightarrow \frac{V_D}{\color{teal}V_T} = {\color{teal}V_T} \ln \left( \frac{ID}{I_S} \right)\\
		\downarrow \\
		V_u = V_i - V_D \cdot \ln \left( \frac{I_D}{I_S} +1 \right)
	\end{gather*}
	Dove $V$ indica la tensione del diodo $V_D$ e dove $V_T$ è pari a $\frac{KT}{q}$.\\\\
	Abbiamo ormai completato lo studio del diodo. \\
	Vogliamo ora campire a cosa può servire il componente elettronico appena ricavato.\\\\
	Consideriamo allora la semplice rete formata da un diodo in rete ad un generatore di tensione $V_I$ e ad una resistenza $R$. (Figura \ref{fig:im-19})\\
	Proviamo allora a ricavare la tensione in uscita $v_u$ a partire dalla differenza di potenziale di ingresso $v_i$.
	\begin{center}
		\begin{figure}[h]
			\centering
			\includegraphics[scale=0.55]{19.png}
			\caption{Un semplice circuito che usa un diodo}
			\label{fig:im-19}
		\end{figure}
	\end{center}
	Detta $V_D$ la caduta di potenziale sul diodo, abbiamo:
	\begin{gather*}
		V_i - V_D - V_u = 0 \\
		\downarrow \\
		V_u = V_i - V_D
	\end{gather*}
	per l’equazione di Kirchhoff alla maglia è:
	\begin{gather*}
		V_u = RI \\
		\downarrow \\
		I_D = \frac{V_u}{R} \\
		\downarrow \\
		{\color{red}V_u = V_i -V_T \cdot \ln \left( \frac{V_u}{R\cdot I_s} +1 \right)}
	\end{gather*}
	per la legge di Ohm. \\
	Introducendo la caratteristica del diodo:
	$$I = I_S \left( e^{\frac{V_D}{V_T}} - 1 \right) = \frac{V_u}{R}$$
	possiamo ricavare:
	$$V_u = RI_S \left( e^{\frac{V_i - V_u}{V_T}} - 1 \right)$$
	Questo risultato è però difficilmente utilizzabile poiché $V_u$ non viene isolato a primo o secondo membro. \\
	Calcoliamo allora la tensione di ingresso.
	$$\frac{V_u}{RI_S} + 1 = e^{\frac{V_i - V_u}{V_T}}$$
	Applichiamo il logaritmo.
	$$V _ { T } \ln \left( { \frac { V _ { u } } { R I _ { S } } } + 1 \right) = { \frac { V _ { i } - V _ { u } } { V _ { T } } }$$
	$$V _ { i } = V _ { u } + V _ { T } \ln \left( { \frac { V _ { u } } { R I _ { S } } } + 1 \right)$$
	Allora siamo riusciti ad esprimere la tensione di ingresso in funzione della tensione di uscita. \\
	Il grafico è riportato in figura \ref{fig:im-20}.
	\begin{center}
		\begin{figure}[h]
			\centering
			\includegraphics[scale=0.55]{20.png}
			\caption{Grafico della tensione di ingresso in funzione della tensione di uscita e grafico della tensione di uscita in funzione di quella di ingresso}
			\label{fig:im-20}
		\end{figure}
	\end{center}
	Siccome la relazione è biunivoca, ci basta ora \textbf{ribaltare il grafico} rispetto alla diagonale del primo e terzo quadrante per ottenere il grafico di $V_u$ in funzione di $V_i$.\\\\
	Dal grafico notiamo che per qualunque valore di tensione in ingresso troviamo in uscita un valore molto simile (anche se leggermente al di sotto) mentre per qualunque valore negativo otteniamo all’uscita un valore molto prossimo allo zero.\\\\
	Supponiamo ora di applicare un segnale variabile nel tempo come una sinusoide. \\
	Sfruttando, anche solo graficamente, la caratteristica del diodo notiamo che per tutti i valori positivi di $v_i$, otteniamo in uscita un valore molto prossimo a quello di $v_i$ mentre per i valori negativi di $v_i$ il diodo appiattisce l’uscita a valori praticamente nulli.
	Allora il diodo è in grado di discriminare le fasi a valori positivi dalle fasi a valori negativi di un segnale tempo variante. \\
	Una prima, evidente, utilità è l’alterazione del valor medio della curva: infatti se in ingresso si ha una sinusoide a valor medio nullo, in uscita avremo un segnale a valor medio positivo. \\
	Un’applicazione estremamente diffusa del diodo è la conversione della tensione da alternata (sinusoidale) ad una tensione continua. \\
	Questo tipo di circuito è detto \textbf{circuito raddrizzatore}.\\\\
	L’efficienza di questa soluzione è alquanto bassa poiché \textbf{annulla metà del segnale}.
	\begin{center}
		\begin{figure}[h]
			\centering
			\includegraphics[scale=0.55]{21.png}
			\caption{Uscita (in rosso) di un raddrizzatore a semionda quando si ha in ingresso (in blu) una sinusoide}
			\label{fig:im-21}
		\end{figure}
	\end{center}
	Scambiamo ora di posizione la resistenza ed il diodo (Figura \ref{fig:im-22}).
	\begin{center}
		\begin{figure}[h]
			\centering
			\includegraphics[scale=0.55]{22.png}
			\caption{Circuito limitatore realizzato con un diodo}
			\label{fig:im-22}
		\end{figure}
	\end{center}
	La tensione di uscita è ora ai capi del diodo. \\
	Applichiamo nuovamente Kirchoff alla maglia.
	$$V_i - RI - V_u = 0 \implies V_u = V_i - RI$$
	$$I = I_S \left( e^{\frac{V_u}{V_T}} - 1 \right)$$
	$$V_u = V_i - RI_S \left( e^{\frac{V_u}{V_T}} - 1 \right)$$
	Anche in questo caso l’equazione è di difficile utilizzo.\\
	 Ripetiamo il procedimento visto nel caso precedente ed otteniamo:
	$$V_i = V_u + RI_S \left( e^{\frac{V_u}{V_T}} - 1 \right)$$
	che ribaltato rispetto alla diagonale porta ad avere il grafico riportato in figura \ref{fig:im-23}.
	\begin{center}
		\begin{figure}[h]
			\centering
			\includegraphics[scale=0.55]{23.png}
			\caption{Uscita di un limitatore}
			\label{fig:im-23}
		\end{figure}
	\end{center}
	Un circuito di questo tipo è un \textbf{circuito limitatore} che viene usato per eliminare tutti i valori di ingresso che eccedono un certo valore. \\
	Questo tipo di circuito è utile quando è necessario proteggere delle componenti da scariche elettrostatiche ad alta tensione.\\\\
	Già da questi semplici esempi si nota che non è sempre possibile ottenere una soluzione analitica esatta dell’equazione.\\
	Possiamo però approssimare la soluzione con metodi di tipo numerico approssimando l’equazione non lineare ad equazioni lineari (gli strumenti elettronici sfruttano questo sistema in alcuni casi). \\
	Una seconda possibilità è quella di risolvere l’equazione l’equazione stessa (e cioè semplificando il modello) in modo da potere ottenere un modello risolvibile analiticamente.
	\subsection{Modello a Soglia}
	Riprendendo la caratteristica del diodo notiamo che essa è composta da due tratti. \\
	Nel tratto di polarizzazione diretta abbiamo che $\frac{dI}{dV}$ è molto grande; nel tratto di polarizzazione inversa invece $\frac{dI}{dV}$ è molto piccola.\\\\
	La curva non è allora approssimabile ad una retta perché perderemmo il carattere di non linearità del diodo. \\
	Tuttavia possiamo sfruttare due rette: una quasi nulla (a livello $-I_S$) per la corrente per la polarizzazione inversa ed una parallela all’asse della corrente (a $V_\gamma \approx 0,6 \sim 0,8$ V) nella regione di polarizzazione diretta.
	\begin{center}
		\begin{figure}[h]
			\centering
			\includegraphics[scale=0.55]{24.png}
			\caption{Caratteristica $V \rightarrow I$ di un diodo nei modelli esponenziale (in rosso) e a soglia (in blu)}
			\label{fig:im-24}
		\end{figure}
	\end{center}
	Allora il modello si semplifica enormemente. \\
	Abbiamo infatti:
	\begin{gather*}
		\begin{aligned}
			\text{Polarizzazione Inversa} \\
			\text{Polarizzazione Diretta}
		\end{aligned}
		\
		\begin{cases}
			I = 0 & \text{per } V < V_\gamma \\ 
			V = V_\gamma & \text{per } I > 0
		\end{cases}
		\qquad
		\begin{aligned}
			\text{Diodo OFF} \\
			\text{Diodo ON}
		\end{aligned}
	\end{gather*}
	Questo modello è detto modello a soglia del diodo poiché discrimina il comportamento del diodo in base ad un valore di soglia.\\\\
	L’approssimazione ci consente di utilizzare due equazioni lineari al posto di una non lineare. \\
	A livello grafico possiamo notare che il modello a soglia fornisce approssimazioni sempre migliori per valori sempre più lontani dal punto di soglia.
	\subsection{Applicazione del Modello a Soglia}
	\subsubsection{Circuito Raddrizzatore}
	Vediamo ora di applicare il modello a soglia al circuito raddrizzatore di Figura 7.7.
	\begin{center}
		\begin{figure}[h]
			\centering
			\includegraphics[scale=0.55]{25.png}
			\caption{Raddrizzatore a semionda}
			\label{fig:im-25}
		\end{figure}
	\end{center}
	Dall’equazione di maglia di Kirchhoff e dalla legge di Ohm abbiamo:
	$$V_i - V - V_u = 0$$
	$$V_u = RI$$
	Incontriamo ora il primo problema del nostro modello approssimato: avendo due regioni di funzionamento dobbiamo risolvere due volte il sistema e controllare che i risultati così ottenuti siano compatibili con le ipotesi di funzionamento formulate\footnote{In generale avendo $n$ diodi, si dovranno analizzare e verificare $2^n$ regioni di funzionamento. Il limite di questo approccio è duplice: bisogna identificare in quali ambiti esso è applicabile e, per grandi circuiti, il numero di calcoli cresce in modo esponenziale.}.\\\\
	Il modello a soglia impone che:
	$$\begin{cases} I = 0 & \text{per } V < V_\gamma \\ 
		V = V_\gamma & \text{per } I > 0 \end{cases}$$
	In questo caso è (Diodo OFF):
	$$\begin{cases} I = 0 \longrightarrow V_u = \underbrace{0}_{R \cdot 0} \\
		V < V_\gamma \ V - V_i - V_u = 0 \end{cases} \implies V_i < V_\gamma$$
	ed il modello risulta corretto.\\\\
	In questo caso è (Diodo ON):
	$$\begin{cases} I > 0 \longrightarrow V_u > 0 \\ 
		V < V_\gamma \longrightarrow \underbrace{V_u = V_i - V_\gamma}_{\textbf{Pendenza 1: Retta traslata}} \end{cases} \implies V_i > V_\gamma$$
	ed il modello risulta corretto.\\\\
	Vediamo il grafico in figura \ref{fig:im-26}
	\begin{center}
		\begin{figure}[h]
			\centering
			\includegraphics[scale=0.55]{26.png}
			\caption{Caratteristica $V \rightarrow I$ di un raddrizzatore}
			\label{fig:im-26}
		\end{figure}
	\end{center}
	\subsubsection{Circuito Limitatore Positivo}
	Prendiamo ora in considerazione il circuito in figura \ref{fig:im-27}
	\begin{center}
		\begin{figure}[h]
			\centering
			\includegraphics[scale=0.55]{27.png}
			\caption{Circuito limitatore positivo}
			\label{fig:im-27}
		\end{figure}
	\end{center}
	Ancora una volta abbiamo:
	$$V_i - RI - V_u = 0$$
	$$V_u = V_D + E$$
	In questo caso abbiamo (Diodo OFF):
	$$\begin{cases} I = 0 \\ 
		I = \frac{V_i - V_u}{R} \text{(Dall'eq. di Kirc qui sopra)}\end{cases} \implies \underbrace{V_i = V_u}_{\text{Bisettrice della prima parte}}$$
	$$\begin{cases} V_D < V_\gamma \\ 
		V_D = V_u - E \end{cases} \implies V_u - E < V_\gamma$$
	da cui ricaviamo:
	$$V_i - E < V_\gamma \implies V_i < E + V_\gamma$$
	Dunque il modello è corretto.\\\\
	In questo caso abbiamo (Diodo ON):
	$$\begin{cases} V_D = V_\gamma \longrightarrow V_u = V_\gamma + E \\ 
		I > 0 \longrightarrow V_i - V_u > 0 \longrightarrow V_i > V_u \end{cases} \implies V_i > E + V_\gamma$$
	Dunque il modello funziona.\\\\
	Vediamo il grafico in figura \ref{fig:im-28}
	\begin{center}
		\begin{figure}[h]
			\centering
			\includegraphics[scale=0.55]{28.png}
			\caption{Caratteristica $V \rightarrow I$ di un limitatore positivo}
			\label{fig:im-28}
		\end{figure}
	\end{center}
	Questo circuito può servire a limitare i picchi di tensione positiva.\\\\
	Si noti che variando il valore di $E$ è possibile aumentare o diminuire arbitrariamente il valore della limitazione.
	\subsubsection{Circuito Limitatore Negativo}
	Prendiamo ora in considerazione il circuito in figura \ref{fig:im-29} (Utilizziamo lo stesso metodo del precedente).
	\begin{center}
		\begin{figure}[h]
			\centering
			\includegraphics[scale=0.55]{29.png}
			\caption{Circuito limitatore negativo}
			\label{fig:im-29}
		\end{figure}
	\end{center}
	Allora per l’equazione di maglia di Kirchhoff è:
	$$V_i + RI = V_u$$
	$$V_u = -E - V_D$$
	In questo caso abbiamo (Diodo OFF):
	$$\begin{cases} I = 0 \longrightarrow V_u = V_i \\ 
		V_D < V_\gamma \\ 
		V_D = -E - V_i \end{cases} \implies -E - V_u < V_\gamma \implies V_u > -E - V_\gamma$$
	In questo caso abbiamo (Diodo ON):
	$$\begin{cases} V_D = V_\gamma \longrightarrow V_u = -(E + V_\gamma) \\ 
		I > 0 \longrightarrow I = \frac{-V_i + V_u}{R} \longrightarrow V_i < V_u \end{cases} \quad \implies V_i < -(E + V_\gamma)$$
	Dunque il modello funziona.
	Vediamo il grafico in figura \ref{fig:im-30}
	\begin{center}
		\begin{figure}[h]
			\centering
			\includegraphics[scale=0.55]{30.png}
			\caption{Caratteristica $V \rightarrow I$ di un limitatore negativo}
			\label{fig:im-30}
		\end{figure}
	\end{center}
	\subsection{Limitatore Doppio}
	FIno ad ora abbiamo incominciato l’analisi di circuiti con componenti non linerari. \\
	Per effettuare questa analisi abbiamo visto due modelli: uno, molto accurato ma spesso difficilmente risolvibile, che utilizza un’equazione esponenziale ed uno, più semplice da risolvere ma meno preciso, che fa uso di due equazioni lineari. \\
	Il secondo modello è chiamato modello regionale ed è riassumibile come:
	$$\begin{cases} I = 0 \quad V_D < V_\gamma \\ 
		I > 0 \quad V_D = V_\gamma \end{cases}$$
	La semplicità viene però pagata con un aumento esponenziale delle equazioni da risolvere: con $n$ diodi dobbiamo infatti considerare $2^n$ regioni (non tutte necessariamente fisicamente possibili).\\\\
	Analizziamo ora il circuito in figra \ref{fig:im-31}\\\\
	Notiamo subito che questo circuito prevede quattro regioni di funzionamento. \\
	Prima di procedere però, dalla configurazione del circuito, ricaviamo:
	$$V_i - V_u = RI$$
	$$I = I_1 - I_2$$
	$$V_u = V_{D1} + E_2 \Longrightarrow V_{D1} = V_u - E_1$$
	$$V_u = -V_{D2} - E_2 \Longrightarrow V_{D2} = -V_u - E_2$$
	Per ipotesi, imponiamo che $E_1$ ed $E_2$ siano positivi.
	\begin{center}
		\begin{figure}[h]
			\centering
			\includegraphics[scale=0.55]{31.png}
			\caption{Circuito per la realizzazione di un limitatore doppio}
			\label{fig:im-31}
		\end{figure}
	\end{center}
	\subsubsection{D1 OFF e D2 OFF}
	Formuliamo ora l’ipotesi che sia D1 sia D2 siano spenti:
	$$\begin{cases} I_1 = 0 \\ 
		I_2 = 0 \end{cases} \implies I = 0 \implies \boxed{V_u = V_i}$$
	$$\begin{cases} V_{D1} < V_\gamma \\ 
		V_{D2} < V_\gamma \end{cases} \implies 
		\begin{cases} V_u - E_1 < V_\gamma \\ 
			-V_u - E_2 < V_\gamma \end{cases} \implies 
			\begin{cases} V_i < E_1 + V_\gamma \\ 
				V_i > -V_\gamma - E_2 \end{cases}$$
	\subsubsection{D1 ON e D2 OFF}
	Consideriamo ora l’ipotesi in cui sia acceso solo D1. \\
	Abbiamo:
	$$\begin{cases} I_1 > 0 \\ 
		I_2 = 0 \end{cases} \implies \underbrace{I > 0}_{I = I_1 - \underbrace{I_2}_{=0} > 0} \implies \frac{V_i - V_u}{R} > 0 \implies V_i > V_u \implies V_i > E_1 + V_\gamma$$
	che è complementare a quanto ottenuto in precedenza. \\
	Controllaimo ora che venga rispettata l’ipotesi che D2 sia spento:
	$$V_{D2} < V_\gamma \implies -(E_2 + V_u) < V_\gamma \implies -(E_2 + E_1 + V_\gamma) < V_\gamma$$
	che risulta essere sempre verificata. \\
	Quindi se il diodo D1 è acceso, allora il diodo D2 non può essere acceso. \\
	Allora possiamo evitare di considerare il caso in cui D1 e D2 siano entrambi accesi.
	\subsubsection{D1 ON e D2 ON}
	Supponiamo comunque, per verifica, che entrambi i diodi siano accesi.\\
	In questo caso è:
	$$\begin{cases} V_{D1} = V_\gamma \implies V_u - E_1 = V_\gamma \implies V_u = E_1 + V_\gamma > 0 \\ 
		V_{D2} = V_\gamma \implies -V_u - E_2 = V_\gamma \implies V_u = -(E_2 + V_\gamma) < 0 \end{cases}$$
	che è \textbf{assurdo} poiché la tensione di uscita non può essere contemporaneamente positiva e negativa.
	\subsubsection{D1 OFF e D2 ON}
	Per finire esaminiamo il caso duale a quello visto nel caso D1 ON e D2 OFF se solo il diodo D2 è acceso risulta essere:
	$$\begin{cases} I_1 = 0 \\ 
		I_2 > 0 \end{cases} \implies I < 0 \implies \frac{V_i - V_u}{R} < 0 \implies V_i < V_u \implies V_i < -(E_2 + V_\gamma)$$
	$$V_{D2} = V_\gamma \implies -V_u - E_2 = V_\gamma \implies V_u = -(E_2 + V_\gamma)$$
	\subsubsection{Grafico e Considerazioni}
	La caratteristica è riportata in figura \ref{fig:im-32} e risulta essere la caratteristica di un limitatore doppio.
	\begin{center}
		\begin{figure}[h]
			\centering
			\includegraphics[scale=0.55]{32.png}
			\caption{Caratteristica di un limitatore doppio}
			\label{fig:im-32}
		\end{figure}
	\end{center}
	Si tratta cioè di un circuito nel quale si hanno tre comportamenti:
	\begin{enumerate}
		\item La tensione di uscita è pari alla tensione di ingresso se questa è compresa nell’intervallo $-(E_2 + V_\gamma) ; (E_1 + V_\gamma)$
		\item La tensione di uscita è limitata al valore $(E_1 + V_\gamma)$ se la tensione di ingresso eccede tale valore
		\item La tensione di uscita è limitata al valore $-(E_2 + V_\gamma)$ se la tensione di ingresso eccede tale valore
	\end{enumerate}
	Questo tipo di circuito risulta molto utile quando è necessario proteggere componenti sensibili da scariche elettriche ad alta intensità.
	\subsection{Rivelatore di Massimo}
	Esaminiamo ora il circuito di Figura \ref{fig:im-33}.
	\begin{center}
		\begin{figure}[h]
			\centering
			\includegraphics[scale=0.55]{33.png}
			\caption{Circuito per la realizzazione di un rivelatore di massimo}
			\label{fig:im-33}
		\end{figure}
	\end{center}
	Grafico in 3 variabili ({\color{red}*DA RIVEDERE!!})\\\\
	Ancora una volta avremo quattro regioni di funzionamento. \\
	Dal circuito possiamo però ricavare preliminarmente:
	$$V_u = RI$$
	$$I = I_1 + I_2$$
	$$V_{D1} = V_1 - V_u$$
	$$V_{D2} = V_2 - V_u$$
	\subsubsection{D1 OFF D2 OFF}
	Abbiamo:
	$$\begin{cases} I_1 = 0 \\ I_2 = 0 \end{cases} \implies I = 0 \implies V_u = 0$$
	che risulta essere verificata per:
	$$\begin{cases} V_{D1} < V_\gamma \\ V_{D2} < V_\gamma \end{cases} \implies \begin{cases} V_1 - V_u < V_\gamma \\ V_2 - V_u < V_\gamma \end{cases} \implies \begin{cases} V_1 < V_\gamma \\ V_2 < V_\gamma \end{cases}$$
	\subsubsection{D1 ON D2 OFF}
	\label{subsub:ril-max-on-off}
	In questo caso risulta essere:
	$$V_{D1} = V_\gamma \implies V_{D1} - V_u = V_\gamma \implies V_u = V_1 - V_\gamma$$
	Le condizioni per cui ciò è verificato sono le seguenti:
	$$\begin{cases} I_1 > 0 \\ I_2 = 0 \end{cases} \implies I > 0 \implies V_u > 0 \implies V_1 > 0$$
	Ed allora ricaviamo:
	$$V_1 > V_\gamma$$
	Ma dobbiamo anche considerare che D2 deve essere spento:
	$$V_{D2} < V_\gamma \implies V_2 - V_u < V_\gamma \implies V_2 < V_u + V_\gamma \implies V_2 < V_1$$
	Allora D2 risulta essere spento se $V_2 < V_1$.
	\subsubsection{D1 ON D2 ON}
	In questa situazione abbiamo:
	$$\begin{cases} V_{D1} = V_\gamma \\ V_{D2} = V_\gamma \end{cases} \implies \begin{cases} V_u = V_1 - V_\gamma \\ V_u = V_2 - V_\gamma \end{cases}$$
	Allora le due espressioni per la tensione di uscita devono portare allo stesso risultato. \\
	Dunque sarà (SSE):
	$$V_1 = V_2$$
	$$\begin{cases} I_1 > 0 \\ I_2 > 0 \end{cases} \implies I > 0 \implies V_u > 0 \implies \begin{cases} V_1 > V_\gamma \\ V_2 > V_\gamma \end{cases}$$
	$$V_u = MAX \left\{ 0; V_1 - V_\gamma; V_2 - V_\gamma \right\}$$
	\subsubsection{D1 OFF D2 ON}
	La situazione è simmettrica rispetto a quanto visto nella sotto-sezione \ref{subsub:ril-max-on-off}. \\
	In questo caso risulta essere:
	$$V_u = V_2 - V_\gamma \implies V_2 > V_\gamma$$
	quando si ha:
	$$V_1 < V_2$$
	\subsubsection{Considerazioni}
	Questo circuito permette di riportare in uscita, a meno di una costante, il segnale di ingresso maggiore tra i due possibili.\\
	Possiamo quindi scrivere che:
	$$V_u = \max \left\{ 0; (V_1 - V_\gamma); (V_2 - V_\gamma) \right\}$$
	Questo circuito è detto \textbf{rivelatore di massimo}.
	\subsubsection{Utilizzo del rivelatore di massimo}
	Il rivelatore di massimo è molto utile quando dobbiamo discriminare i valori di un segnale binario digitale. \\
	Ipotizziamo di assegnare il valore $V_L \approx 0$ al simbolo 0 e $V_H \approx 5$ al simbolo 1.\\\\
	La Tabella \ref{tab:table-3} ne riassume il comportamento.
	\begin{table}[h]
		\centering
		\begin{tabular}{|c|c|c|c|c|c|c|c|}
			\hline
			$V_1$ & $V_2$ & $D_1$ & $D_2$ & $V_u$ & $x_1$ & $x_2$ & $y$ \\ \hline
			\multicolumn{8}{|c|}{$\quad$} \\ [-0.8em] \hline
			$V_L$ & $V_L$ & OFF & OFF & $V_L$ & 0 & 0 & 0 \\ \hline
			$V_L$ & $V_H$ & OFF & ON  & $V_H$ & 0 & 1 & 1 \\ \hline
			$V_H$ & $V_L$ & ON  & OFF & $V_H$ & 1 & 0 & 1 \\ \hline
			$V_H$ & $V_H$ & ON  & ON  & $V_H$ & 1 & 1 & 1 \\ \hline
		\end{tabular}
		\caption{Comportamento digitale di un rivelatore di massimo e sua versione in logica booleana}
		\label{tab:table-3}
	\end{table}
	\\Dalla tabella notiamo che il rivelatore di massimo realizza una \textbf{porta OR a due ingressi}. \\
	Aumentando arbitrariamente il numero di diodi possiamo aumentare arbitrariamente il numero di ingressi.\\\\
	Questa tecnologia è estremamente semplice ed economica ma non garantisce prestazioni elevate.
	\subsection{Rilevatore di Minimo}
	Analizziamo adesso il circuito di Figura \ref{fig:im-34}.
	\begin{center}
		\begin{figure}[h]
			\centering
			\includegraphics[scale=0.55]{34.png}
			\caption{Circuito per la realizzazione di un rivelatore di minimo}
			\label{fig:im-34}
		\end{figure}
	\end{center}
	Nuovamente dall’analisi circuitale ricaviamo:
	$$V_u = E - RI$$
	$$I = I_1 + I_2$$
	$$V_{D1} = V_u - V_1$$
	$$V_{D2} = V_u - V_2$$
	\subsubsection{D1 OFF D2 OFF}
	Abbiamo:
	$$\begin{cases} I_1 = 0 \\ I_2 = 0 \end{cases} \implies I = 0 \implies V_u = E$$
	che risulta essere verificata per:
	$$\begin{cases}V_{D1}<V_\gamma\\V_{D2}<V_\gamma&\end{cases}\Longrightarrow\begin{cases}V_1-V_u<V_\gamma\\V_2-V_u<V_\gamma&\end{cases}\Longrightarrow\begin{cases}V_1>E-V_\gamma\\V_2>E-V_\gamma&\end{cases}$$
	\subsubsection{D1 ON D2 OFF}
	\label{subsub:ril-min-on-off}
	In questo caso risulta essere:
	$$V_{D1} = V_{\gamma} \implies V_u - V_1 = V_{\gamma} \implies V_u = V_1 + V_{\gamma}$$
	Le condizioni per cui ciò è verificato sono le seguenti:
	$$\begin{cases}I_1>0\\I_2=0&\end{cases}\implies I>0\implies\frac{E-V_u}{R}>0\implies E>V_u\implies V_1<E-V_\gamma$$
	Ma dobbiamo anche considerare che D2 deve essere spento:
	$$V_{D2} < V_{\gamma} \implies V_u - V_2 < V_{\gamma} \implies V_2 > V_u - V_{\gamma} \implies V_2 > V_1$$
	Allora D2 risulta essere spento se $V_2 > V_1$.
	\subsubsection{D1 ON D2 ON}
	In questa situazione abbiamo:
	$$\begin{cases}V_{D1}=V_\gamma\\V_{D2}=V_\gamma\end{cases}\Longrightarrow\begin{cases}V_u=V_1+V_\gamma\\V_u=V_2+V_\gamma&\end{cases}$$
	Allora le due espressioni per la tensione di uscita devono portare allo stesso risultato. 
	Dunque sarà:
	$$V_1 = V_2$$
	$$\begin{cases}I_1>0\\I_2>0&\end{cases}\implies I>0\implies V_u>0\implies\begin{cases}V_1<E-V_\gamma\\V_2<E-V_\gamma&\end{cases}$$
	\subsubsection{D1 OFF D2 ON}
	La situazione è simmettrica rispetto a quanto visto nella sottosezione \ref{subsub:ril-min-on-off}.\\
	In questo caso risulta essere:
	$$V_u = V_2 + V_{\gamma} \implies V_2 < E - V_{\gamma}$$
	quando si ha:
	$$V_1 > V_2$$
	\subsubsection{Considerazioni}
	Questo circuito (detto rivelatore di minimo) ha la funzione di riportare in uscita il valore più piccolo tra quelli fornitigli in ingresso a meno di una costante $V_{\gamma}$.
	$$V_u = \min \left\{ E; (V_1 + V_{\gamma}); (V_2 + V_{\gamma}) \right\}$$
	\subsubsection{Utilizzo del rilevatore di minimo}
	Interpretiamo il rivelatore di minimo in senso digitale. \\
	Di nuovo ipotizziamo di assegnare il valore $V_L \approx 0$ al simbolo 0 e $V_H \approx 5$ al simbolo 1.\\\\
	La Tabella \ref{tab:table-4} ne riassume il comportamento.
	\begin{table}[h]
		\centering
		\begin{tabular}{|c|c|c|c|c|c|c|c|}
			\hline
			$V_1$ & $V_2$ & $D_1$ & $D_2$ & $V_u$ & $x_1$ & $x_2$ & $y$ \\
			\hline
			\multicolumn{8}{|c|}{$\quad$}\\[-0.8em] \hline% Small empty gap row
			$V_L$ & $V_L$ & ON  & ON  & $V_L$ & 0 & 0 & 0 \\ \hline
			$V_L$ & $V_H$ & ON  & OFF & $V_L$ & 0 & 1 & 0 \\ \hline
			$V_H$ & $V_L$ & OFF & ON  & $V_L$ & 1 & 0 & 0 \\ \hline
			$V_H$ & $V_H$ & OFF & OFF & $V_H$ & 1 & 1 & 1 \\ 
			\hline
		\end{tabular}
		\caption{Comportamento digitale di un rivelatore di massimo e sua versione in logica booleana}
		\label{tab:table-4}
	\end{table}
	\\Dalla tabella notiamo che il rivelatore di massimo realizza una \textbf{porta AND a due ingressi}.\\
	Aumentando arbitrariamente il numero di diodi possiamo aumentare arbitrariamente il numero di ingressi.\\\\
	Analogamente alla porta OR realizzata con i diodi anche questa porta AND è estremamente semplice ed economica ma non garantisce prestazioni elevate.
	\subsection{Raddrizzatore a doppia semionda}
	Analizziamo il circuito in figura \ref{fig:im-35}
	\begin{center}
		\begin{figure}[h]
			\centering
			\includegraphics[scale=0.55]{35.png}
			\caption{Raddrizzatore a doppia semionda}
			\label{fig:im-35}
		\end{figure}
	\end{center}
	A prima vista, trattandosi di un circuito con quattro diodi, \textbf{dovremo analizzare sedici regioni di funzionamento differenti}. \\
	La Tabella 9.1 mostra le sedici possibili regioni. \\
	Sono indicate con una “X” le regioni fisicamente non ammesse.
	\begin{table}[h]
		\centering
		\begin{tabular}{|c|c|c|c|c|c|c}
			\cline{3-7}
			\multicolumn{2}{c|}{} & OFF & OFF & ON & ON & \multicolumn{1}{c|}{\textbf{D3}} \\ \cline{3-7}
			\multicolumn{2}{c|}{} & OFF & ON & OFF & ON & \multicolumn{1}{c|}{\textbf{D4}} \\ \hline
			OFF & OFF & \ref{subsub:raddr-sopp-sem-off-off-off-off} & X & X & X & \multicolumn{1}{c}{} \\ \cline{1-6}
			OFF & ON  & X & X & \ref{subsub:raddr-sopp-sem-off-on-on-off} & X & \multicolumn{1}{c}{} \\ \cline{1-6}
			ON  & OFF & X & \ref{subsub:raddr-sopp-sem-on-off-off-on} & X & X & \multicolumn{1}{c}{} \\ \cline{1-6}
			ON  & ON  & X & X & X & X & \multicolumn{1}{c}{} \\ \cline{1-6}
			\textbf{D1} & \textbf{D2} & \multicolumn{5}{c}{} \\ \cline{1-2}
		\end{tabular}
		\caption{Tabella delle regioni di funzionamento del raddrizzatore a doppia semionda. \\ Le “X” indicano le zone fisicamente non possibili}
		\label{tab:table-5}
	\end{table}
	Analizzando preliminarmente il circuito possiamo ricavare che sia (Kirchoff alla maglia e al nodo):
	$$I = I_1 + I_3 = I_4 + I_2$$
	$$V_u = RI$$
	$$V_{D1} + V_{D2} + V_u = 0$$
	$$V_{D3} + V_{D4} + V_u = 0$$
	$$V_i - V_{D1} - V_u -V_{D3} = 0$$
	$$V_i + V_{D2} + V_u +V_{D4} = 0$$
	\subsubsection{D1 ON}
	Ipotizzando che il diodo $D1$ è acceso, otteniamo dalle equazioni di maglia che è:
	$$\begin{cases} I_1 > 0 \\ I_3 > 0 \\ I_4 \geq 0\end{cases} \implies I > 0 \implies V_u > 0$$
	$$V_{D1} - V_\gamma > 0$$
	$$V_{D2} = -(V_u + V_{D1})$$
	e di conseguenza avremo:
	$$V_{D2} < 0 < V_\gamma$$
	Ma allora quando il diodo $D1$ è acceso il diodo $D2$ sarà necessariamente spento.\\\\
	Sempre dalle equazioni di maglia notiamo quando $D1$ è acceso, sarà anche acceso il diodo $D4$. \\
	Infatti sarà:
	$$I_4 = I - I_2 \implies V_{D4} - V_\gamma > 0$$
	$$V_{D3} = -(V_u + V_{D4}) < 0$$
	Ed allora anche quando il diodo $D4$ risulta acceso il diodo $D3$ non può essere acceso.
	\subsubsection{D2 ON}
	Se il diodo $D2$ è acceso avremo:
	$$\begin{cases} I_2 > 0 \\ I_4 \geq 0 \end{cases} \implies I > 0 \implies V_u > 0$$
	$$V_{D2} = V_\gamma > 0$$
	$$V_{D1} = -(V_u + V_{D2}) < 0$$
	Ma allora se il diodo $D2$ è acceso, il diodo $D1$ deve essere spento e quindi:
	$$I_1 = 0$$
	Pertanto, siccome è:
	$$I = I_3 + I_1$$
	avremo che la corrente è fornita dal diodo $D3$ che dovrà dunque essere acceso.\\\\
	Ancora, sempre dalle equazioni di maglia, notiamo che il diodo non sono ragionevoli neppure le condizioni di accensione di un solo diodo.\\\\
	Ma allora \textbf{le regioni da analizzare si sono ridotte da sedici a tre}.
	\subsubsection{D1, D2, D3 e D4 OFF}
	\label{subsub:raddr-sopp-sem-off-off-off-off}
	In questa ipotesi la corrente è evidentemente nulla.
	$$I_1 = I_2 = I_3 = I_4 = 0 \implies I = 0$$
	Allora per la legge di Ohm, la tensione di uscita è nulla.
	$$V_u = 0$$
	Questa condizione è accettabile quando si ha:
	$$\begin{cases} V_{D1} < V_\gamma \\ V_{D2} < V_\gamma \\ V_{D3} < V_\gamma \\ V_{D4} < V_\gamma \end{cases}$$
	Usando le equazioni di maglia per i diodi $D1$ e $D4$ abbiamo:
	$$V_i - V_{D1} - V_u - V_{D4} = 0 \implies V_{D1} + V_{D4} = V_i$$
	da cui segue che deve essere:
	$$V_i < 2V_\gamma$$
	Procedendo nello stesso modo per $D2$ e $D3$, avremo:
	$$V_i + V_{D2} + V_{D3} + V_u = 0 \implies V_{D2} + V_{D3} = -V_i$$
	e quindi:
	$$V_i > -2V_\gamma$$
	Allora l’ambito di validità di questa regione è:
	$$-2V_\gamma < V_i < 2V_\gamma$$
	\subsubsection{D1 ON, D2 OFF, D3 OFF e D4 ON}
	\label{subsub:raddr-sopp-sem-on-off-off-on}
	Le ipotesi significano che la corrente è:
	$$I_1 = I_4 = I > 0$$
	Ma è anche:
	$$V_{D1} = V_{D4} = V_\gamma$$
	Sfruttiamo l’equazione di Kirchhoff vista nel sottosezione \ref{subsub:raddr-sopp-sem-off-off-off-off}:
	$$V_i - V_{D1} - V_u - V_{D4} \implies V_u = V_i - 2V_\gamma$$
	Questa condizione è verificata quando si ha:
	$$V_u > 0$$
	che, in relazione alla tensione di ingresso, è:
	$$V_i > 2V_\gamma$$
	\subsubsection{D1 OFF, D2 ON, D3 ON e D4 OFF}
	\label{subsub:raddr-sopp-sem-off-on-on-off}
	È la condizione duale a quelle della sottosezione \ref{subsub:raddr-sopp-sem-on-off-off-on} abbiamo:
	$$I_3 = I_2 = I > 0$$
	$$\begin{cases} V_{D2} = V_\gamma \\ V_{D3} = V_\gamma \end{cases}$$
	Sfruttiamo nuovamente l’equazione di maglia di Kirchhoff vista nella sottosezione \ref{subsub:raddr-sopp-sem-off-off-off-off}
	$$V_i + V_{D2} + V_{D3} + V_u \Longrightarrow V_u = -(V_i + V_{D2} + V_{D3}) \Longrightarrow V_u = -V_i - 2V_\gamma$$
	Queste ipotesi sono valide quando si ha:
	$$V_i < -2V_\gamma$$
	\subsubsection{Grafico e considerazioni}
	Il grafico della caratteristica è riportato in Figura \ref{fig:image5}
	\begin{figure}[h]
		
		\begin{subfigure}{0.5\textwidth}
			\includegraphics[width=0.9\linewidth]{36.png} 
			\caption{Caratteristica del raddrizzatore a doppia semionda}
			\label{fig:im-36}
		\end{subfigure}
		\begin{subfigure}{0.5\textwidth}
			\includegraphics[width=0.9\linewidth]{37.png}
			\caption{}
			\label{fig:im-37}
		\end{subfigure}
		
		\caption{}
		\label{fig:image5}
	\end{figure}
	Questo circuito si comporta come un raddrizzatore a semionda per i valori positivi di $V_i$ che eccedono la soglia $2V_\gamma$. \\
	I valori di ingresso negativi che eccedono il valore $-2V_\gamma$ vengono riportati in positivo (vengono cioè “raddrizzati”).\\\\
	Per via di questo comportamento il circuito prende il nome di raddrizzatore a doppia semionda.\\\\
	Il raddrizzatore a doppia semionda è molto più efficiente di un raddrizzatore a semionda poiché, pur consentendo di ottenere un segnale a valor medio non nullo da un segnale a valor medio nullo, non elimina metà del segnale e, raddrizzandolo, porta ad un valor medio più elevato a parità di ingresso.\\\\
	La configurazione esaminata è detta ponte raddrizzatore di diodi ed è spesso disegnata come in Figura \ref{fig:im-38}
	\begin{center}
		\begin{figure}[h]
			\centering
			\includegraphics[scale=0.55]{38.png}
			\caption{Rappresentazione grafica classica del raddrizzatore a doppia semionda}
			\label{fig:im-38}
		\end{figure}
	\end{center}
	\subsection{Convertitore alternata continua}
	\label{sub:ac-dc}
	I circuito raddrizzatori visti fino a questo momento, pur raddrizzando il segnale e portando un valor medio non nullo all’uscita, riportano comunque segnali periodici non costanti. \\
	Molti dispositivi domestici funzionano però a corrente continua. \\
	Dobbiamo trovare dunque un modo per estrarre la componente continua del segnale.\\\\
	Questo risultato si può ottenere facilmente aggiungendo un filtro passa basso che blocchi le componenti ad alta frequenza del segnale. \\
	Il più semplice filtro passa basso possibile è un condensatore che si comporta come un circuito aperto per frequenze nulle e come un corto circuito per frequenze elevate.\\\\
	Analizziamo ora cosa succede aggiungendo un filtro passa basso al raddrizzatore a singola semionda (per semplicità). \\
	Lo schema è riportato in Figura \ref{fig:im-39}
	\begin{center}
		\begin{figure}[h]
			\centering
			\includegraphics[scale=0.55]{39.png}
			\caption{L’aggiunta del condensatore di capacità $C$ consente di ridurre la variabilità del segnale di uscita}
			\label{fig:im-39}
		\end{figure}
	\end{center}
	L’aggiunta del condensatore di capacità $C$ modifica la caratteristica. \\
	Per le relazioni costitutive di resistenza e condensatore e per la legge di Kirchhoff abbiamo ora:
	$$I_R = \frac{V_u}{R}$$
	$$I_C = C \cdot \frac{dV_u}{dt}$$
	$$I = I_R + I_C \implies I = \frac{V_u}{R} + C \cdot \frac{dV_u}{dt}$$
	$$V_i - V_D - V_u = 0$$
	Ipotizziamo di avere in ingresso:
	$$V_i = V_m \sin(\omega t)$$
	con $V_m \gg V_\gamma$ (che è tipicamente verificata).
	\subsubsection{Diodo acceso}
	In questo caso è:
	$$V_D = V_\gamma \implies V_u = V_m \sin(\omega t)$$
	$$I > 0 \implies \frac{V_u}{R} + C \cdot \frac{dV_u}{dt} > 0$$
	Ma, per come è definito l’ingresso, il diodo sarà certamente attivo nell’intervallo $\left[0; \frac{T}{4}\right]$ dove $T$ è il periodo. \\
	Nell’intervallo $\left[\frac{T}{2}; \frac{3T}{4}\right]$ il diodo sarà invece necessariamente spento. \\
	Tuttavia nell’intervallo $\left(\frac{T}{4}; \frac{T}{2}\right)$ in cui si avrà $I = 0$ e cioè un punto, $t_{\text{OFF}}$ nel quale il diodo si spegne.
	\subsubsection{Spegnimento del diodo}
	Per calcolare $t_{\text{OFF}}$ imponiamo che sia:
	$$\frac{V_u}{R} + C \cdot \frac{dV_u}{dt} = 0$$
	$$-\frac{V_u}{R} = C \cdot \frac{dV_u}{dt}$$
	$$-\frac{V_m \sin(\omega t)}{R} = C \cdot \frac{d}{dt} [V_m \sin(\omega t)]$$
	$$-\frac{V_m}{R} \cdot \sin(\omega t) = V_m C \omega \cdot \cos(\omega t)$$
	$$\frac{\sin(\omega t)}{\cos(\omega t)} = -\omega RC$$
	$$\tan(\omega t_{\text{OFF}}) = -\omega RC \implies \omega t_{\text{OFF}} = -\arctan(\omega RC) \pm k\pi$$
	Allora abbiamo ricavato l’istante in cui il diodo si spegne. \\
	Ricordando che la tangente tende a $\pm\infty$ per $\alpha \rightarrow \pm\frac{\pi}{2}$ ed è periodica su un intervallo di durata $T = \pi$ è chiaro che il punto di spegnimento del diodo si troverà nell’angolo compreso tra $\frac{\pi}{2}$ e $\pi$.\\\\
	Riportato sul grafico di Figura \ref{fig:im-40} vediamo che il punto di spegnimento è tanto più vicino al tempo $\frac{T}{4}$ quanto più è elevato il valore di $RC$.
	Da $t_{\text{OFF}}$ in avanti il diodo risulta spento e la corrente sul diodo è nulla.
	\begin{center}
		\begin{figure}[h]
			\centering
			\includegraphics[scale=0.55]{40.png}
			\caption{Quando il diodo si spegne il condensatore inizia la sua fase di scarica}
			\label{fig:im-40}
		\end{figure}
	\end{center}
	$\omega_{t_{OFF}}$ è detto \textbf{Angolo di } {\color{red}* DA SISTEMARE!!!}
	\subsubsection{Diodo spento}
	Quando il diodo è spento abbiamo:
	$$I = 0$$
	$$\frac{V_u}{R} + C \cdot \frac{dV_u}{dt} = 0$$
	Ora però l’ingresso è nullo ed il circuito è semplificato come in Figura \ref{fig:im-41}
	\begin{center}
		\begin{figure}[h]
			\centering
			\includegraphics[scale=0.55]{41.png}
			\caption{Quando il diodo si spegne il circuito si semplifica ad un parallelo tra il condensatore e la resistenza}
			\label{fig:im-41}
		\end{figure}
	\end{center}
	Risolvendo l’equazione differenziale abbiamo:
	$$\frac{1}{RC} \int_{t_{\text{OFF}}}^t dt = \int_{V_u(t_{\text{OFF}})}^{V_u(t)} \frac{1}{V_u} dV_u$$
	$$\frac{t - t_{\text{OFF}}}{RC} = \ln \frac{V_u(t)}{V_u(t_{\text{OFF}})}$$
	$$V_u(t) = V_u(t_{\text{OFF}}) e^{\frac{t - t_{\text{OFF}}}{RC}}$$
	Allora l’uscita decresce con una curva esponenziale tanto più velocemente quanto più è piccolo il valore di $RC$.\\\\
	Questa condizione è verificata per tutti i valori per cui si ha:
	$$V_u > V_i$$
	Allora quando incontriamo nuovamente la sinusoide di ingresso torniamo ad utilizzare la relazione usata per il diodo acceso e quindi torniamo ad avere:
	$$V_u = V_i$$
	\subsection{Problemi del diodo}
	Dalle simulazioni notiamo che, dimensionando opportunamente le costanti, possiamo utilizzare il convertitore esaminato nella sezione \ref{sub:ac-dc} per estrarre una funzione di uscita che si appoggi sul massimo del segnale di ingresso. \\
	Questa applicazione è molto utile quando si tratta di demodulare un segnale modulato in ampiezza.\\\\
	Fino a questo momento abbiamo applicato segnali tempo varianti all’ingresso di un circuito costituito da uno o più diodi. \\
	Aumentando la frequenza (senza aumentare le costanti $R$ e $C$) ci accorgiamo però di una anomalia non prevista dal modello.\\\\
	Il grafico di Figura \ref{fig:im-42} mostra come il diodo modifichi la sua caratteristica $V_i \to V_u$ all’aumentare della frequenza (il diodo fatica sempre più a spegnersi). \\
	Per una frequenza troppo elevata si arriva al limite per il quale il diodo non si spegne mai e continua a seguire sempre la curva di ingresso.
	\begin{center}
		\begin{figure}[h]
			\centering
			\includegraphics[scale=0.55]{42.png}
			\caption{Per frequenze elevate il diodo a giunzione $pn$ fatica sempre più a spegnersi \\ Per frequenze sufficientemente alte, il diodo resta sempre acceso}
			\label{fig:im-42}
		\end{figure}
	\end{center}
	Il modello visto fino ad ora ha funzionamento per due motivi:
	\begin{enumerate}
		\item Abbiamo usato il modello statico del diodo alla contare la variabile ( $t$ )
		\item Abbiamo sempre usato frequenze basse
	\end{enumerate}
	Tuttavia il diodo ha bisogno di un tempo non nullo per potersi spegnere. \\
	Allora per un periodo $T \to 0$, il diodo non avrà il tempo fisico di spegnersi.\\\\
	Nel grafico di Figura \ref{fig:im-43} l’area azzurra rappresenta la concentrazione di lacune che si spostano nel passaggio da polarizzazione diretta a polarizzazione inversa.
	\begin{center}
		\begin{figure}[h]
			\centering
			\includegraphics[scale=0.55]{43.png}
			\caption{Differenti concentrazioni di lacune per tensioni differenti: in rosso la situazione di equilibrio; in verde la condizione di polarizzazione inversa $V < 0$; in blu la condizione di polarizzazione diretta $V > 0$. \\ L’area azzurra è l’area necessaria allo spegnimento del diodo.}
			\label{fig:im-43}
		\end{figure}
	\end{center}
	Allora la carica non si può spostare in un tempo nullo (sarebbe necessaria un’energia infinita) e quindi abbiamo un tempo di reazione (effetti reattivi associati alla giunzione $p_n$) del diodo che sarà a sua volta non nullo.\\\\
	Il modello usato fino ad ora (\textbf{modello statico}) è utilizzabile solo per piccoli scostamenti dall’equilibrio e per segnali lentamente variabili nel tempo.
	\subsection{Modello Dinamico}
	Iniziamo ora lo studio di un modello migliore del diodo che tenga conto della variazione nel tempo. \\
	All’inizio dello studio del diodo a giunzione $p_n$ avevamo calcolato che la densità di carica nella regione svuotata fosse:
	$$Q_n = -Q_p = -\underbrace{qN_A}_{\rho} w_p$$
	dove:
	$$w_p = \frac{1}{N_A} \sqrt{\frac{2\varepsilon\psi_{B0}}{q\left(\frac{1}{N_A} + \frac{1}{N_D}\right)}} \implies Q_p = \sqrt{\frac{2\varepsilon q\psi_{B_0}}{\left(\frac{1}{N_A} + \frac{1}{N_D}\right)}}$$
	In ipotesi di piccolo scostamente avevamo inoltre visto che:
	$$\psi_B = \psi_{B_0} - V$$
	Allora possiamo ricavare la carica dinamica del diodo in funzione della tensione:
	$$Q(V) = \sqrt{\frac{2\varepsilon q(\psi_B - V)}{\left(\frac{1}{N_A} + \frac{1}{N_D}\right)}}$$
	che, nuovamente, è una funziune non lineare.\\\\
	Si noti che per $V = 0$ il valore della carica $Q_0 = Q(0)$ non passa per l’origine. \\
	La capacità del diodo è data dalla derivata della curva $Q(V)$
	$$DQ(V) = \frac{d}{dV}Q(V)$$
	e quindi, per comodità, useremo la curva $Q - Q_0$ in modo da ottenere $Q(0) = 0$ che non modifica la derivata.\\\\
	In questa condizione il diodo si comporta come un condensatore di capacità variabile (\textbf{diodo varicap}) ed era molto utilizzato nei primi televisori analogici a colori.
	Applicando una tensione $V > \psi_{B_0}$ abbiamo invece una soluzione assurda: la base della radice è negativa. \\
	In realtà, per $V = \psi_{B_0}$ abbiamo che il campo elettrico e la barriera di potenziale sono entrambi nulli. \\
	Ma se la barriera è nulla, la regione svuotata non esiste e di conseguenza non ha più senso parlare di piccolo scostamento.\\\\
	Allora la pendenza della curva della carica del diodo è detta capacità di svuotamento. \\
	In particolare la pendenza in $V = 0$ è detta capacità di giunzione all’equilibrio.\\\\
	Siccome la regione svuotata non esiste, dobbiamo considerare, per la carica, un secondo contributo dovuto alla carica dei portatori mobili. \\
	Si può dimostrare che il secondo contorbuto di carica dipende esponenzialmente dalla tensione applicata (cosa prevedibile in quanto la corrente cresce esponenzialmente in polarizzazione diretta). \\
	Questo secondo conturbuto è trascurabile per la polarizzazione inversa mentre è importante per la polarizzazione diretta.
	$$\begin{cases} J \propto \left(e^{\frac{V}{V_T}} - 1\right) \\ J \propto \frac{dn}{dx} \end{cases} \implies \frac{dn}{dx} \propto \left(e^{\frac{V}{V_T}} - 1\right) \implies n \propto \left(e^{\frac{V}{V_T}} - 1\right)$$
	Possiamo allora ottenere un modello un poco approssimato ma comunque sufficientemente preciso prendendo entrambe le componenti dominanti della carica nel diodo.
	$$\begin{cases} Q(V) = \sqrt{\psi_{B0} - V} & V < 0 \\ Q(V) = Q_S \left(e^{\frac{V}{V_T}} - 1\right) & V \geq 0 \end{cases}$$
	Tale modello è raffigurato in figura \ref{fig:im-44}.
	\begin{center}
		\begin{figure}[h]
			\centering
			\includegraphics[scale=0.55]{44.png}
			\caption{Grafico del modello dinamico del diodo. In rosso la funzione $Q$ per $V < 0$; in blu la funzione $Q - Q_0$; in verde la funzione $Q$ per $V \geq 0$}
			\label{fig:im-44}
		\end{figure}
	\end{center}
	Riprendendo il modello statico possiamo sfruttare il fatto che sia:
	$$\begin{cases} I_0 = I_D \left(e^{\frac{V}{V_T}} - 1\right) \\ Q = Q_S \left(e^{\frac{V}{V_T}} - 1\right) \end{cases} \implies \frac{Q}{I_0} = \frac{Q_S \left(e^{\frac{V}{V_T}} - 1\right)}{I_D \left(e^{\frac{V}{V_T}} - 1\right)} = \frac{Q_S}{I_D} = \tau$$
	possiamo scrivere la carica in funzione della corrente od il contrario:
	$$Q = \tau I_D \iff I_D = \frac{Q}{\tau}$$
	Il valore $\tau$ è un tempo ed è il\textbf{tempo di vita medio dei portatori}. \\
	È cioè il tempo medio intercorso tra il momento in cui un portatore $p$ diffonde in zona $n$ ed il momento in cui esso si ricombina.\\\\
	Allora abbiamo ricavato un modello dinamico fortemente non lineare del diodo. \\
	Nuovamente introduciamo una semplificazione simile a quella vista per la corrente.
	$$\begin{cases}V<V_\gamma&\begin{cases}I_D=0\\Q\approx0\end{cases} \\ V=V_\gamma&\begin{cases}I_D>0\\Q>0\end{cases}\end{cases}$$
	Questo modello approssimato per la carica è meno preciso rispetto al modello statico visto per la carica. \\
	Il diodo dinamico è rappresentato come in Figura \ref{fig:im-45}.
	\begin{center}
		\begin{figure}[h]
			\centering
			\includegraphics[scale=0.55]{45.png}
			\caption{Simbolo del modello dinamico del diodo}
			\label{fig:im-45}
		\end{figure}
	\end{center}
	La corrente al terminale di uscita del diodo è data da:
	$$I = I_{D} + \frac{dQ}{dt}$$
	Allora per frequenze basse possiamo riutilizzare le relazioni studiate per il caso statico (per tali frequenze $\frac{dQ}{dt}$ è infatti trascurabile). \\
	Questo diodo si può modellare sia con il controllo di corrente:
	$$I = I_{D} \left( e^{\frac{V}{V_{T}}} - 1 \right)$$
	sia con il modello a controllo di carica:
	$$I = \frac{Q}{\tau} = \frac{Qs \left( e^{\frac{V}{V_{T}}} - 1 \right)}{\tau}$$
	\subsection{Analisi DInamica}
\end{document}